\documentclass[11pt]{article}

\usepackage{amsmath,amssymb,amsthm}
\usepackage{geometry}
\usepackage{hyperref}
\usepackage{macros}

\geometry{margin=1in}

\title{
Structural Classification of Undecidable Problems\\
via Non-Middle Fallacy and Dynamic Normalization
}

\author{
Masaomi Chiba
}

\date{2025-12-14}

\begin{document}

\maketitle

\begin{abstract}
Many central open problems in modern mathematics exhibit a common
structural difficulty: they attempt to determine static truth values
for phenomena that are inherently dynamic, inter-structural, or
self-referential.
In this paper, we introduce the notion of a \emph{Non-Middle Fallacy} (NMF),
a structural obstruction arising when a static formal system attempts
to enforce absolute equality across incompatible logical regimes.

We provide a unified classification of major undecidable or controversial
problems---including the ABC conjecture, the Riemann hypothesis,
P versus NP, the Continuum Hypothesis, and Godel-type incompleteness---
according to distinct types of NMF.
We then propose a dynamic normalization framework, denoted $\mathbf{mweq}$,
which replaces static equality with a stability-based equivalence.
This framework does not claim to produce classical proofs, but rather
to normalize structural tension and classify the logical status of
statements beyond the theorem/hypothesis dichotomy.

Finally, we introduce a refined taxonomy of mathematical truth
(Nomological Truth, Praxis Truth, and Quasi-Theorem),
intended to clarify the epistemic status of results that are
structurally stable yet classically undecidable.
\end{abstract}

\section{Introduction}

\subsection{Motivation}

A remarkable feature of modern mathematics is that many of its most
important problems resist resolution not because of technical
difficulty alone, but because of deeper structural constraints.
Examples include the ABC conjecture, the Riemann hypothesis,
the P versus NP problem, and the Continuum Hypothesis.
Despite decades of effort, these problems remain undecided within
standard formal frameworks.

At the same time, foundational results such as Godel's incompleteness
theorems and Cohen's independence results demonstrate that undecidability
is not an anomaly but an intrinsic feature of sufficiently expressive
formal systems.
This raises a natural question:

\begin{quote}
Are these difficulties isolated phenomena, or manifestations of a
common structural pattern?
\end{quote}

This paper argues for the latter.

\subsection{Overview of the Approach}

We introduce the concept of a \emph{Non-Middle Fallacy} (NMF), a structural
configuration in which a static formal system attempts to enforce
absolute equality or decidability across domains that require
dynamic mediation.
Rather than treating undecidability as a failure, we treat it as a
diagnostic signal of structural mismatch.

Our approach proceeds in three stages:

\begin{enumerate}
  \item Define NMFs as structural obstructions arising from static
        formalization.
  \item Classify major undecidable or controversial problems according
        to distinct types of NMF.
  \item Introduce a dynamic normalization framework, $\mathbf{mweq}$,
        that replaces static equality with stability-based equivalence.
\end{enumerate}

The goal is not to refute classical mathematics, but to extend its
descriptive power by clarifying the structural status of problems
that lie beyond classical proof.

\section{Non-Middle Fallacy (NMF)}

\subsection{Definition}

We define a \emph{Non-Middle Fallacy} (NMF) as follows.

\medskip
\noindent
\textbf{Definition 1 (Non-Middle Fallacy).}
An NMF occurs when a static formal system attempts to enforce a binary,
absolute decision (true/false, equal/unequal) on a phenomenon whose
logical structure is inherently dynamic, inter-structural, or
self-referential, thereby producing undecidability, inconsistency,
or persistent controversy.
\medskip

NMFs do not indicate an error in reasoning, but rather a mismatch
between the expressive mode of the formal system and the structure
of the problem.

\subsection{Structural Origins of NMF}

We identify three recurring structural sources of NMFs.

\begin{enumerate}
  \item \textbf{Dynamic vs Static Mismatch:}
        Finite static descriptions are applied to infinite or
        process-oriented phenomena.
  \item \textbf{Inter-Structural Comparison:}
        Distinct logical domains are compared via static equality.
  \item \textbf{Boundary and Self-Reference:}
        A system attempts to fully describe or decide properties of
        its own expressive boundary.
\end{enumerate}

These sources give rise to distinct classes of NMF, which we analyze
in the next sections.


\section{Classification of Non-Middle Fallacy}

In this section, we present a unified classification of Non-Middle Fallacy
based on their structural origin.
Although individual problems differ greatly in content, their logical
pathologies exhibit strong structural similarities.

\subsection{Inter-Structural Comparison NMF}

\subsubsection{Description}

An \emph{Inter-Structural Comparison NMF} arises when two logically
distinct structures are compared using static equality.
The comparison itself is well-defined syntactically, but lacks a
stable semantic interpretation.

Typical examples involve comparisons between:
\begin{itemize}
  \item additive and multiplicative structures,
  \item discrete and continuous domains,
  \item computational processes and verification procedures.
\end{itemize}

\subsubsection{Representative Problems}

\begin{center}
\begin{tabular}{|l|l|p{7cm}|}
\hline
Problem & Status & Structural Description \\
\hline
ABC conjecture & Undecided / disputed &
Static comparison between additive relation $a+b=c$
and multiplicative radical $rad(abc)$. \\
\hline
Riemann hypothesis & Undecided &
Static comparison between discrete prime distribution
and analytic continuation of zeta zeros. \\
\hline
Goldbach conjecture & Undecided &
Additive decomposition constrained by multiplicative
(prime) structure. \\
\hline
P vs NP & Undecided &
Equality between dynamic computation time
and static verification time. \\
\hline
IUT treatment of ABC & Controversial &
Comparison across meta-arithmetic universes with
restricted identification rules. \\
\hline
\end{tabular}
\end{center}

\subsubsection{Structural Analysis}

In all such cases, the difficulty does not lie in the absence of local
arguments, but in the lack of a coherent static notion of equality
that bridges the two domains.
Attempts to resolve these problems classically tend to increase
structural complexity rather than reduce logical tension.

\subsection{Boundary and Self-Reference NMF}

\subsubsection{Description}

A \emph{Boundary/Self-Reference NMF} occurs when a formal system attempts
to decide properties that inherently refer to its own expressive limits.
This type of NMF is characterized by unavoidable meta-level escalation.

\subsubsection{Representative Problems}

\begin{center}
\begin{tabular}{|l|l|p{7cm}|}
\hline
Problem & Status & Structural Description \\
\hline
Godel incompleteness & Proven undecidable &
Finite axiomatic system attempts to decide its own
consistency. \\
\hline
Chaitin incompleteness & Proven undecidable &
Finite description attempts to capture irreducible
algorithmic information. \\
\hline
Russell paradox & Resolved by stratification &
Unrestricted self-membership within a static definition. \\
\hline
\end{tabular}
\end{center}

\subsubsection{Structural Analysis}

These results are often interpreted negatively as limitations.
From the NMF perspective, they instead indicate that static closure
is the wrong objective.
The failure is not logical inconsistency, but the absence of a
dynamic stabilization mechanism.

\subsection{Hierarchical Inclusion NMF}

\subsubsection{Description}

A \emph{Hierarchical Inclusion NMF} arises when an infinite generative
hierarchy is forced into a static axiomatic determination.
This is typical in set-theoretic questions concerning infinity.

\subsubsection{Representative Problems}

\begin{center}
\begin{tabular}{|l|l|p{7cm}|}
\hline
Problem & Status & Structural Description \\
\hline
Continuum Hypothesis & Independent of ZFC &
Static decision of infinite cardinal hierarchy. \\
\hline
Large cardinal axioms & Independent of ZFC &
Existence claims beyond fixed axiomatic reach. \\
\hline
\end{tabular}
\end{center}

\subsubsection{Structural Analysis}

Independence results show that these questions do not admit a
privileged static resolution.
Any answer requires extending the universe of discourse, confirming
their classification as NMF phenomena rather than incomplete proofs.

\subsection{Unified Perspective}

Despite surface differences, all NMF types share a common pattern:
a static binary decision is imposed on a structure that demands
dynamic mediation.
This observation motivates the search for a normalization framework
that replaces static equality with a stability-based criterion.


\section{Dynamic Normalization via mweq}

This section introduces the core technical contribution of this paper:
the replacement of static equality by a dynamic stability operator,
denoted by $mweq$.
This operator provides a normalization mechanism that resolves
Non-Middle Fallacy without increasing structural complexity.

\subsection{Limitations of Static Equality}

Classical equality is defined as a timeless, context-independent
relation.
While effective in homogeneous domains, it becomes unstable when
applied across structurally distinct systems.

Formally, expressions of the form
\[
A = B
\]
implicitly assume:
\begin{itemize}
  \item identical generative rules,
  \item shared invariants,
  \item identical semantic resolution.
\end{itemize}

In the presence of NMF, these assumptions fail, producing either
undecidability or explosive complexity.

\subsection{Definition of the mweq Operator}

We define $mweq$ as a relation between structures that captures
dynamic convergence rather than static identity.

\begin{definition}[Middle-Way Equivalence]
Let $S_1$ and $S_2$ be two structures equipped with transformations
$F$ and $G$.
We write
\[
S_1 \; mweq \; S_2
\]
if there exists a finite sequence of transformations such that the
structural distortion measure $E$ satisfies
\[
E(F(G(x))) \rightarrow E_{min}
\]
for all admissible $x$.
\end{definition}

Here $E$ is not required to vanish.
Stability is defined by convergence to a minimal admissible value,
not by exact coincidence.

\subsection{Finite Stabilization Principle}

A central principle of $mweq$-based mathematics is the exclusion of
static infinity.

\begin{principle}[Finite Stabilization]
All valid constructions must admit a finite stabilization process.
Infinite objects are replaced by arbitrarily extensible but finite
procedures whose stability can be evaluated dynamically.
\end{principle}

This principle eliminates the need for the axiom of infinity and
replaces it with iterative convergence criteria.

\subsection{Resolution of Inter-Structural Comparison NMF}

Given two structures $S_1$ and $S_2$, classical mathematics asks whether
they are equal or not.
The $mweq$ framework instead asks whether their interaction stabilizes.

\subsubsection{Functorial Coherence}

Let
\[
f : S_1 \rightarrow S_2, \quad g : S_2 \rightarrow S_1
\]
be structure-preserving maps.
The condition
\[
g \circ f \; mweq \; id_{S_1}
\]
expresses dynamic coherence rather than strict invertibility.

This replaces static comparison with a convergence condition that can
be verified algorithmically.

\subsection{Resolution of Boundary and Self-Reference NMF}

Boundary-related NMF are resolved by abandoning static completeness.

\subsubsection{Meta-Level Stabilization}

Let $C$ be a formal system and $M$ a meta-system.
Instead of requiring $C$ to decide its own consistency, we introduce
a reflective transformation
\[
R : C \rightarrow M
\]
such that instability in $C$ triggers corrective normalization in $M$.

Consistency is thus maintained dynamically as a fixed point of
iteration, rather than a statically provable property.

\subsection{Resolution of Hierarchical Inclusion NMF}

Hierarchical questions such as the continuum hypothesis rely on
static determination of infinite strata.

Under $mweq$, such questions are reformulated as follows:
\begin{quote}
Does a given generative process stabilize under finite extension?
\end{quote}

If stabilization occurs, the hierarchy is admissible.
If not, the question itself is declared structurally invalid rather
than undecidable.

\subsection{Summary}

The $mweq$ operator provides a low-complexity alternative to classical
meta-theoretic escalation.
By replacing equality with dynamic stability, it resolves all major
classes of Non-Middle Fallacy while preserving computational and
conceptual tractability.


\section{New Classes of Mathematical Truth}

This section introduces a reclassification of mathematical truth based
on dynamic coherence rather than static provability.
The classical dichotomy of theorem versus conjecture is shown to be
insufficient in the presence of Non-Middle Fallacy.

\subsection{Limitations of the Theorem--Conjecture Dichotomy}

In classical mathematics, a statement is assigned one of two statuses:

\begin{itemize}
  \item \textbf{Theorem}: provable within a fixed axiom system by a
  finite sequence of deductions.
  \item \textbf{Conjecture}: not yet provable or undecidable within the
  given system.
\end{itemize}

This classification presupposes that all meaningful statements admit
static truth values.
However, NMF demonstrates that many foundational questions fail not
because of missing techniques, but because the question itself is
structurally ill-posed.

\subsection{Dynamic Truth and Coherence Levels}

Under the $mweq$ framework, truth is evaluated by the degree to which a
statement maintains dynamic stability across structural transformations.

We therefore introduce three classes of truth, ordered by increasing
strength of coherence.

\subsection{Nomological Truth}

\begin{definition}[Nomological Truth]
A statement is called a \emph{nomological truth} if it expresses a
structural constraint required for universal dynamic stability.
Such a statement resolves all associated NMF and remains coherent
under arbitrary admissible transformations.
\end{definition}

Nomological truths are not contingent on a particular axiom system.
They correspond to invariants of the middle-way normalization process.

\paragraph{Example.}
The reformulated ABC principle, expressed as a bound on structural
distortion between additive and multiplicative constructions, is a
candidate for nomological truth.

\subsection{Practical Truth}

\begin{definition}[Practical Truth]
A statement is called a \emph{practical truth} if it guarantees dynamic
stability under finite operational procedures, even if universal
coherence is not established.
\end{definition}

Practical truths are characterized by effective convergence rather
than absolute invariance.

\paragraph{Example.}
Monte Carlo methods exemplify practical truth: finite sampling does not
establish exact equality, but reliably converges toward stable outcomes.

\subsection{Quasi-Theorem}

\begin{definition}[Quasi-Theorem]
A statement is called a \emph{quasi-theorem} if it is internally
consistent and provable within a newly constructed formal universe,
but lacks verified coherence with external or classical systems.
\end{definition}

This class legitimizes innovative frameworks without prematurely
forcing classical equivalence.

\paragraph{Example.}
Results derived within inter-universal Teichmueller theory are best
understood as quasi-theorems until functorial coherence with classical
arithmetic is demonstrated.

\subsection{Structural Classification of Major Problems}

The following table summarizes the proposed classification.

\begin{center}
\begin{tabular}{|l|l|l|}
\hline
Problem & NMF Type & Truth Class \\
\hline
ABC conjecture & Inter-structural & Nomological \\
Riemann hypothesis & Inter-structural & Nomological \\
P vs NP & Dynamic/static & Nomological \\
Monte Carlo methods & Dynamic/static & Practical \\
IUT results & Inter-structural & Quasi-theorem \\
Continuum hypothesis & Hierarchical & Invalid question \\
\hline
\end{tabular}
\end{center}

\subsection{Consequences for Mathematical Practice}

This reclassification shifts the goal of mathematics from absolute
closure to controlled coherence.
Innovation is no longer blocked by premature demands for static
equivalence, while foundational confusion is reduced by explicit
recognition of NMF.

\section{Conclusion}

This paper has argued that many of the most persistent problems in
mathematics arise not from technical difficulty, but from a shared
structural defect, the Non-Middle Fallacy (NMF).
Static formal systems attempt to force dynamic, relational, or
self-referential phenomena into fixed logical identities, thereby
producing undecidability, paradox, or endless controversy.

\subsection{Summary of Results}

The main contributions of this work are as follows.

\begin{enumerate}
  \item A unified definition of NMF was introduced, explaining
  undecidability, incompleteness, and foundational conflict as a single
  structural phenomenon.
  \item The operator $mweq$ was proposed as a dynamic replacement for
  static equality, allowing comparisons based on convergence and
  stability rather than absolute identity.
  \item A normalization framework based on finite operations and
  dynamic middle-way coherence was shown to dissolve major classes of
  NMF.
  \item Classical open problems were reclassified according to their
  structural nature rather than their proof status.
  \item Three new truth classes were introduced: nomological truth,
  practical truth, and quasi-theorem.
\end{enumerate}

These results suggest that the limits identified by Goedel and later
formal results do not merely constrain provability, but signal the need
for a change in the underlying notion of mathematical truth.

\subsection{Implications for Foundations}

Under the proposed framework, mathematics is no longer required to be
self-contained in a static sense.
Instead, it is understood as a dynamically stabilized system whose
coherence is maintained through iterative normalization.

This perspective resolves several long-standing tensions:

\begin{itemize}
  \item Incompleteness becomes a structural boundary condition rather
  than a failure.
  \item Independence results are reinterpreted as indicators of invalid
  static questions.
  \item Highly complex theories are evaluated by coherence rather than
  by forced equivalence with classical formalisms.
\end{itemize}

In particular, inter-universal constructions may be meaningful without
serving as classical proofs, provided their internal dynamics are
stable.

\subsection{Outlook}

The framework developed here opens multiple directions for future
research.

\begin{itemize}
  \item In number theory, major conjectures may be reformulated as
  nomological truths expressing unavoidable structural constraints.
  \item In logic and computer science, dynamic normalization provides a
  new lens on complexity, verification, and computation.
  \item In physics and information theory, middle-way coherence offers
  a foundation for understanding stability without singularities.
\end{itemize}

More broadly, this work suggests that the role of mathematics is
shifting from the pursuit of absolute identities to the cultivation of
stable relations.
The adoption of dynamic middle-way principles may therefore mark a
transition comparable in significance to the axiomatization movement
of the twentieth century.

\subsection{Final Remark}

The Non-Middle Fallacy is not an obstacle to be eliminated, but a signal
that a static formulation has exceeded its domain of validity.
By embracing dynamic coherence through $mweq$, mathematics can proceed
beyond its classical limits without abandoning rigor.

This shift does not weaken mathematics.
It clarifies its scope, strengthens its foundations, and restores
contact between formal reasoning and the structures it seeks to
understand.

\end{document}
